% !TeX program = xelatex
% UTF-8 transcription of the handwritten source; medical claims are preserved.
% pic01.jpg = 20261013_035145904.jpg (page 1)
% pic02.jpg = 20261013_035152146.jpg (page 2)
% pic03.jpg = 20261013_035156909.jpg (page 3)
% Compile directly with XeLaTeX, or input into a document using xepersian.
\ifdefined\settextfont
  \def\OsteomyelitisStandalone{0}
\else
  \def\OsteomyelitisStandalone{1}
  \documentclass[12pt,a4paper]{article}
  \usepackage[margin=2cm]{geometry}
  \usepackage{amsmath,amssymb}
  \usepackage{xepersian}
  \IfFontExistsTF{Amiri}{\settextfont{Amiri}}{\settextfont{FreeSerif}}
  \setlatintextfont{FreeSerif}
\fi
\providecommand{\en}[1]{\lr{#1}}
\newcommand{\osteopage}[1]{\clearpage\section*{استئومیلیت --- صفحهٔ #1}}
\ifnum\OsteomyelitisStandalone=1 \begin{document}\fi

\osteopage{۱}
\subsection*{انواع و محل درگیری}
\begin{itemize}
\item شایع‌ترین فرم استئومیلیت در کودکان: \en{Acute Hematogenous}.
\item استئومیلیت تحت حاد: به دنبال تروماهای نافذ موضعی.
\item استئومیلیت مزمن: اگر فرم تحت حاد به‌درستی درمان نشود.
\item در کودکانی که هموگلوبینوپاتی ندارند، استئومیلیت تقریباً همواره در \textbf{متافیز استخوان‌های بلند} رخ می‌دهد.
\item استئومیلیت در کودکان با بیماری سیکل سل: در \textbf{دیافیز استخوان‌های بلند}.
\item در کودکان زیر سن ۱۲--۱۸ ماه، به دلیل وجود شریان‌ها، عفونت از متافیز به اپی‌فیز منتقل می‌شود و \textbf{آرتریت} ایجاد می‌کند.
\end{itemize}

\subsection*{پاتوژن‌ها}
\begin{itemize}
\item شایع‌ترین پاتوژن: \textbf{استاف اورئوس}.
\item دیگر پاتوژن‌ها: استرپ گروه \en{A}، استرپ پنومونیه، \en{TB}، بارتونلا، نایسریا.
\item \textbf{سالمونلا، استاف اورئوس}: استئومیلیت به دنبال هموگلوبینوپاتی و بیماری \en{SSD}.
  % The source abbreviation is SSD; it has not been changed to SCD.
\item گازگرفتگی سگ و گربه: \textbf{پاستورلا}.
\item گازگرفتگی انسان: \textbf{کینگلا}.
\item شایع‌ترین پاتوژن استئومیلیت مزمن: استاف اورئوس.
\end{itemize}

\subsection*{استئومیلیت مولتی‌فوکال عودکنندهٔ مزمن}
\begin{latin}
Chronic recurrent multifocal osteomyelitis (chronic non bacterial)
\end{latin}
\begin{itemize}
\item بیماری \textbf{خودالتهابی}، غیرعفونی، با اپیزودهای راجعه.
\item تب، درد استخوان، یافته‌های تصویربرداری استئومیلیت.
\item محل درگیری: کلاویکل، اسکاپولا، استخوان‌های کوچک دست و پا.
\end{itemize}

\subsection*{سن شایع}
\begin{itemize}
\item سن شایع استئومیلیت: \textbf{۱۲--۳۰ سال}؛ \textbf{۲ برابر خانم‌ها} در \en{M}.
\end{itemize}

\osteopage{۲}
\subsection*{تظاهرات بالینی}
\begin{itemize}
\item درد موضعی، تندرنس نقطه‌ای، گرمی، اریتم، تورم.
\item تب، بی‌اشتهایی، تحریک‌پذیری، بی‌قراری.
\item توانایی تحمل وزن و حرکات فعال و غیرفعال کم می‌شود: \textbf{سودوپارالیز}.
\end{itemize}

\subsection*{آزمایش‌ها}
\begin{itemize}
\item لکوسیتوز: ممکن است باشد یا نباشد.
\item \en{BC}: \textbf{لازم است} ولی اغلب منفی می‌شود.
\item \textbf{\en{\(\mathrm{ESR},\ \mathrm{CRP}\uparrow\)}}.
\end{itemize}

\subsection*{تصویربرداری}
\begin{itemize}
\item گرافی معمولی، تورم موضعی و جابه‌جایی اتصالات (تندینوس) را ظرف ۳ روز نخست نشان می‌دهد.
\item ضایعات استخوانی ۴--۱۴ روز بعد در گرافی دیده می‌شوند.
\item \textbf{\en{MRI} حساس‌ترین تصویربرداری است}.
\end{itemize}

\subsection*{درمان و اندیکاسیون‌های جراحی}
\begin{itemize}
\item علی‌رغم درمان، اگر تا \textbf{۴۸ ساعت} تب رفع نشد: اندیکاسیون قطعی با \textbf{جراحی}.
\item سایر اندیکاسیون‌های جراحی:
  \begin{itemize}
  \item تشکیل آبسه.
  \item درگیری مفصل هیپ.
  \item \en{Sequestrum}.
  \item اسپاینال کورد کامپرشن (اورژانس).
  \end{itemize}
\item طول درمان آنتی‌بیوتیکی: حداقل \textbf{۳--۴ هفته}.
\item بدترین پروگنوز: نوزادان، درگیری مفصل \textbf{شانه و هیپ}.
\item پس از پایین آمدن \en{CRP} و \en{ESR} و بهبود بالینی: ترخیص و ادامهٔ درمان به‌صورت \en{Oral} یا \en{IV} در منزل.
\end{itemize}

\osteopage{۳}
\subsection*{انتخاب آنتی‌بیوتیک}
\subsubsection*{هماتوژن حاد}
\begin{itemize}
\item استاف اورئوس:
  \begin{itemize}
  \item حساس به متی‌سیلین: نافسیلین، اگزاسیلین، سفازولین.
  \item \textbf{مقاوم به متی‌سیلین}: کلیندا، ونکو، تری‌متوپریم--سولفامتوکسازول (\en{TMP-SMX}).
  \end{itemize}
\item استرپ پیوژن (گروه \en{A}): آمپی‌سیلین / آموکسی‌سیلین.
\item \textbf{استرپ پنومونیه}: آمپی‌سیلین / آموکسی‌سیلین، سفتریاکسون، ونکو.
\end{itemize}

\subsubsection*{موضعی تحت حاد}
\begin{itemize}
\item \textbf{سودومونا آئروژینوزا}: سفتازیدیم، پیپراسیلین--تازوباکتام، سیپروفلوکساسین.
\item کینگلا: \textbf{آمپی‌سیلین}، نسل ۳ سفالوسپورین.
\end{itemize}

\ifnum\OsteomyelitisStandalone=1
  \let\OsteomyelitisEnd\enddocument
\else
  \let\OsteomyelitisEnd\relax
\fi
\OsteomyelitisEnd
